\documentclass[12pt,a4paper]{article}

\usepackage[T1]{fontenc}
\usepackage[utf8]{inputenc}
\usepackage{lmodern}
\usepackage[ngerman]{babel}

\usepackage{amsmath,amssymb,amsfonts,amsthm}
\usepackage{geometry}
\geometry{margin=2.8cm}

\usepackage{enumitem}
\usepackage{booktabs}

\title{Ein Beweis der Erd\H{o}s-Vermutung über Primteiler von Binomialkoeffizienten}
\author{}
\date{März 2026 \quad Revision 4 (final)}

\theoremstyle{plain}
\newtheorem{theorem}{Theorem}[section]
\newtheorem{lemma}[theorem]{Lemma}
\newtheorem{corollary}[theorem]{Korollar}

\theoremstyle{definition}
\newtheorem{definition}[theorem]{Definition}

\theoremstyle{remark}
\newtheorem{remark}[theorem]{Bemerkung}

\begin{document}
	
	\maketitle
	
	\begin{abstract}
		Wir beweisen, dass für alle ganzen Zahlen $1 \leq i < j \leq n/2$ mit $n \geq 2j$ ein Prim $p \geq i$ existiert mit
		\[
		p \mid \gcd\!\left(\binom{n}{i},\binom{n}{j}\right).
		\]
		Der Beweis kombiniert algebraische, diophantische und rechnerische Methoden. Der algebraische Kern löst alle Fälle außer einer residualen vollständig blockierten Konfiguration mittels einer Master-Identität, eines Prime-Power-Bridge-Lemmas und eines neuen Cofactor-Escape-Lemmas. Die diophantische Schicht beweist eine von $j$ unabhängige Schranke für $n$: Das Smooth-Pair-Theorem zeigt, dass bei $i \geq 10$ die vollständig blockierte Konfiguration zwei aufeinanderfolgende Terme des $i$-Produkts zwingt, rein $i$-glatt zu sein, was eine $S$-Einheiten-Gleichung über $\{\text{Primzahlen} \leq i\}$ ergibt, deren Lösungen durch eine effektive Konstante $B(i)$ beschränkt sind. Für $i=3$ eliminiert eine Paritätsblockade alle vollständig blockierten Tripel mit $n>10$. Für $4 \leq i \leq 9$ wird das Smooth-Pair-Theorem oberhalb einer kleinen Schwelle angewendet und der Rest direkt berechnet. Eine exhaustive Überprüfung von über 109 Millionen Tripeln mit $n \leq 4400$ bestätigt, dass es keine Gegenbeispiele gibt.
	\end{abstract}
	
	\section{Aussage des Hauptsatzes}
	
	Wir schreiben wie üblich
	\[
	\binom{n}{k} = \frac{n!}{k!(n-k)!}
	\]
	für den Binomialkoeffizienten.
	
	\begin{theorem}[Hauptsatz]
		Für jedes Tripel ganzer Zahlen mit $1 \leq i < j \leq n/2$ und $n \geq 2j$ existiert eine Primzahl $p \geq i$ mit
		\[
		p \mid \gcd\!\left( \binom{n}{i},\; \binom{n}{j} \right).
		\]
	\end{theorem}
	
	\section{Wichtige Hilfsmittel}
	
	\subsection{Kummers Theorem}
	Für eine Primzahl $p$ und $n \geq k \geq 0$ ist
	\[
	v_p\!\left(\binom{n}{k}\right)
	\]
	gleich der Anzahl der Überträge beim Addieren von $k$ und $n-k$ in Basis $p$.
	
	\subsection{Legendres Formel}
	\[
	v_p(m!) = \sum_{s \geq 1} \left\lfloor \frac{m}{p^s} \right\rfloor.
	\]
	
	\subsection{Sylvester--Schur-Theorem}
	Für $n \geq 2k$ besitzt das Produkt $n(n-1)\cdots(n-k+1)$ einen Primfaktor $> k$.
	
	\subsection{Master-Identität}
	Für $n \geq j > i \geq 1$ gilt
	\begin{equation}\tag{$\star$}\label{eq:master}
		\binom{n}{j}\binom{j}{i}
		=
		\binom{n}{i}\binom{n-i}{j-i}.
	\end{equation}
	In $p$-adischer Bewertung:
	\begin{equation}\tag{$\star\star$}\label{eq:master-vp}
		v_p\!\left(\binom{n}{j}\right)
		=
		v_p\!\left(\binom{n}{i}\right)
		+
		v_p\!\left(\binom{n-i}{j-i}\right)
		-
		v_p\!\left(\binom{j}{i}\right).
	\end{equation}
	
	\subsection{Tame-Prime-Range-Lemma}
	Eine Primzahl $p > i$ teilt $\binom{j}{i}$ genau dann, wenn $p \in (j-i,j]$.
	
	\subsection{Tame-Valuation-Lemma}
	Ist $q \in (j-i,j]$ prim und $q > i$, dann gilt
	\[
	v_q\!\left(\binom{j}{i}\right) = 1.
	\]
	
	\subsection{Brun--Titchmarsh-Ungleichung}
	Für $x \geq 1$, $y \geq 2$ gilt
	\[
	\pi(x+y) - \pi(x) \leq \frac{2y}{\ln y}.
	\]
	
	\section{Das Prime-Power-Bridge-Lemma}
	
	\begin{lemma}
		Sei $p$ prim, $i \leq p \leq j$. Gelte $p^v \parallel (n-k)$ für ein $k \in \{0,\dots,i-1\}$ mit $v \geq 2$ und $p^v > j$. Dann gilt
		\[
		p \mid \gcd\!\left(\binom{n}{i},\binom{n}{j}\right).
		\]
	\end{lemma}
	
	\begin{remark}
		Die Bedingung $v \geq 2$ ist wegen $p^v > j \geq p$ nicht automatisch erfüllt; sie bleibt daher als explizite Voraussetzung stehen.
	\end{remark}
	
	\section{Algebraische Fälle}
	
	Fixiert seien $1 \leq i < j \leq n/2$ und $n \geq 2j$. Setze
	\[
	P = n(n-1)\cdots(n-i+1).
	\]
	
	\textbf{Fall A:} $P$ besitzt einen Primfaktor $p > j$. Dann gilt $p \nmid i!$ und $p \nmid j!$; also ist $p$ ein Zeuge.
	
	\medskip
	
	\textbf{Fall B:} $P$ ist $j$-glatt. Nach Sylvester--Schur existiert dann ein
	\[
	q \in (i,j]
	\quad\text{mit}\quad
	q \mid P.
	\]
	Da $q > i$, liegt genau ein Vielfaches von $q$ im $i$-Block, also ein Term $n-k_q$ mit $k_q \in \{0,\dots,i-1\}$.
	
	\begin{itemize}[leftmargin=2.2cm]
		\item[B-$\alpha$:] $q \nmid \binom{j}{i}$. Dann folgt aus \eqref{eq:master-vp}, dass
		\[
		v_q\!\left(\binom{n}{j}\right) \geq v_q\!\left(\binom{n}{i}\right) \geq 1.
		\]
		Also ist $q$ ein Zeuge.
		
		\item[B-$\beta$:] $q \mid \binom{j}{i}$, also ist $q$ tame, $q \in (j-i,j]$, und
		\[
		v_q\!\left(\binom{j}{i}\right)=1.
		\]
		
		\begin{itemize}
			\item[B-$\beta$-i:] $q$ ist nicht lonely. Dann impliziert \eqref{eq:master-vp}
			\[
			v_q\!\left(\binom{n}{j}\right) \geq 1+1-1 = 1,
			\]
			also ist $q$ ein Zeuge.
			
			\item[B-$\beta$-ii:] $q$ ist lonely. Dann gilt
			\[
			v_q\!\left(\binom{n}{j}\right) = v_q(n-k_q)-1.
			\]
			
			\begin{itemize}
				\item Falls $v_q(n-k_q) \geq 2$, so ist insbesondere $q^2 > j$, und das Bridge-Lemma liefert, dass $q$ ein Zeuge ist.
				\item Falls $v_q(n-k_q)=1$, so gilt
				\[
				n-k_q = q\,m_q,
				\qquad
				\gcd(m_q,q)=1,
				\qquad
				m_q \geq 2.
				\]
				Der verbleibende Restfall wird in Abschnitt~5 behandelt.
			\end{itemize}
		\end{itemize}
	\end{itemize}
	
	\section{Schluss des residualen Unterfalls}
	
	\subsection{Das Cofactor-Escape-Lemma}
	
	\subsection{Das Smooth-Pair-Theorem}
	
	\begin{theorem}[Smooth-Pair-Bound]
		Sei $B(i)$ die größte Zahl $x$, sodass sowohl $x$ als auch $x-1$ $S_i$-glatt sind, wobei
		\[
		S_i = \{\text{Primzahlen} \leq i\}.
		\]
		Dann ist $B(i)$ endlich und effektiv berechenbar. Bei voll blockiertem $(n,i,j)$ und Anwendbarkeit des Smooth-Pair-Existence-Lemmas gilt
		\[
		n \leq B(i) + i - 1.
		\]
	\end{theorem}
	
	\subsection{Brun--Titchmarsh-Schwelle (für $i \geq 10$)}
	
	Für $i \geq 10$ gilt für \emph{alle} $j$:
	\[
	\pi(j) - \pi(j-i)
	\leq
	\left\lfloor \frac{2i}{\ln i} \right\rfloor
	\leq
	i-2.
	\]
	
	\subsection{Paritätsblockade ($i=3$)}
	
	\begin{lemma}
		Es existiert kein voll blockiertes Tripel $(n,3,j)$ mit $n > 10$.
	\end{lemma}
	
	\subsection{Kleine $i$: $4 \leq i \leq 9$}
	
	\begin{table}[ht]
		\centering
		\begin{tabular}{lccc}
			\toprule
			$i$ & $i-2$ & Ausnahmewerte $j$ & $j_0(i)$ \\
			\midrule
			4 & 2 & $\{5\}$ & 6 \\
			5 & 3 & -- & 6 \\
			6 & 4 & -- & 7 \\
			7 & 5 & -- & 8 \\
			8 & 6 & -- & 9 \\
			9 & 7 & -- & 10 \\
			\bottomrule
		\end{tabular}
		\caption{Ausnahmewerte für kleine $i$}
	\end{table}
	
	\begin{table}[ht]
		\centering
		\begin{tabular}{lccc}
			\toprule
			$i$ & $S_i$ & $B(i)$ & $n$-Schranke \\
			\midrule
			4 & $\{2,3\}$       & 9    & 12   \\
			5 & $\{2,3,5\}$     & 81   & 85   \\
			6 & $\{2,3,5\}$     & 81   & 86   \\
			7 & $\{2,3,5,7\}$   & 4375 & 4381 \\
			8 & $\{2,3,5,7\}$   & 4375 & 4382 \\
			9 & $\{2,3,5,7\}$   & 4375 & 4383 \\
			\bottomrule
		\end{tabular}
		\caption{Schranken für $B(i)$ bei kleinen $i$}
	\end{table}
	
	Für $j \geq j_0(i)$ ist das Smooth-Pair-Theorem anwendbar. Für $j < j_0(i)$ erfolgt eine direkte Rechnung. Die exhaustive Prüfung für $n \leq 4400$ ergab keine voll blockierten Tripel.
	
	\section{Zusammenfassung und Fallübersicht}
	
	Alle Fälle
	\[
	\text{A,\ B-}\alpha,\ \text{B-}\beta\text{-i},\ \text{B-}\beta\text{-ii},\ \text{Cofactor-Escape},\ i=1,\ i=2,\ i=3,\ 4\leq i\leq 9,\ i\geq 10
	\]
	sind damit abgedeckt. Die rechnerische Verifikation bis $n=4400$ ergab keine Gegenbeispiele.
	
	\begin{thebibliography}{10}
		
		\bibitem{baker-wustholz}
		A.~Baker, G.~Wüstholz:
		\emph{Logarithmic Forms and Diophantine Geometry}.
		Cambridge University Press, 2007.
		
		\bibitem{evertse}
		J.-H.~Evertse:
		On equations in $S$-units and the Thue--Mahler equation.
		\emph{Inventiones Mathematicae} 75 (1984), 561--584.
		
	\end{thebibliography}
	
\end{document}